\chapter{2008年上海卷（秋文）}

\section{填空题}

\begin{problem}
不等式 \(|x - 1| < 1\) 的解集是 \fillinblank{} \fillinblank{}.
\end{problem}

\begin{problem}
若集合 \(A = \{x \mid x \leqslant 2\}\)、\(B = \{x \mid x \geqslant a\}\) 满足 \(A \cap B = \{2\}\)，则实数 \(a = \)  \fillinblank{}.
\end{problem}

\begin{problem}
若复数 z 满足 \( z = \i (2 - z) \) (i 是虚数单位)，则 \( z = \)  \fillinblank{}.
\end{problem}

\begin{problem}
若函数 \( f(x) \) 的反函数为 \( f^{-1}(x)=\log_{2}x \) ，则 \( f(x)= \)  \fillinblank{}.
\end{problem}

\begin{problem}
若向量 \(a, b\) 满足 \(|a| = 1, |b| = 2\)，且 \(a\) 与 \(b\) 的夹角为 \(\frac{\pi}{3}\)，则 \(|a + b| = \)  \fillinblank{}.
\end{problem}

\begin{problem}
若直线 \(ax - y + 1 = 0\) 经过抛物线 \(y^{2} = 4x\) 的焦点，则实数 \(a = \)  \fillinblank{}.
\end{problem}

\begin{problem}
若 \(z\) 是实系数方程 \(x^{2} + 2x + p = 0\) 的一个虚根，且 \(|z| = 2\) ，则 \(p = \)  \fillinblank{}.
\end{problem}

\begin{problem}
在平面直角坐标系中，从五个点: \(A(0,0), B(2,0), C(1,1), D(0,2), E(2,2)\) 中任取三个，这三点能构成三角形的概率是 \fillinblank{}. (结果用分数表示)
\end{problem}

\begin{problem}
若函数 \( f(x) = (x + a)(bx + 2a) \) （常数 \( a, b \in \R \)）是偶函数，且它的值域为 \((- \infty, 4]\)，则该函数的解析式 \( f(x) = \)  \fillinblank{}.
\end{problem}

\begin{problem}
已知总体的各个体的值由小到大依次为 2, 3, 3, 7, a, b, 12, 13, 17, 18, 3, 20，且总体的中位数为 10.5. 若要使该总体的方差最小，则 a, b 的取值分别是 \fillinblank{} \fillinblank{}.
\end{problem}

\begin{problem}
在平面直角坐标系中，点 \(A\)、\(B\)、\(C\) 的坐标分别为 \((0,1)\)、\((4,2)\)、\((2,6)\). 如果 \(P(x,y)\) 是 \(\triangle ABC\) 围成的区域（含边界）上的点，那么当 \(w = xy\) 取到最大值时，点 \(P\) 的坐标是 \fillinblank{} \fillinblank{}.
\end{problem}

\section{单选题}

\begin{problem}
设 \(P\) 是椭圆 \(\frac{x^2}{25} + \frac{y^2}{16} = 1\) 上的点. 若 \(F_1, F_2\) 是椭圆的两个焦点，则 \(\left|PF_1\right| + \left|PF_2\right| = \frac{25}{16} \neq\)（\quad）

\choices
    {4}
    {5}
    {8}
    {10}
\end{problem}

\begin{problem}
给定空间中的直线 \(l\) 及平面 \(\alpha\)，条件“直线 \(l\) 与平面 \(\alpha\) 内无数条直线都垂直”是“直线 \(l\) 与平面 \(\alpha\) 垂直”的（\quad）

\choices
    {充分非必要条件}
    {必要非充分条件}
    {充要条件}
    {既非充分又非必要条件}
\end{problem}

\begin{problem}
若数列 \(\{a_{n}\}\) 是首项为 1，公比为 \(a - \frac{3}{2}\) 的无穷等比数列，且 \(\{a_{n}\}\) 各项的和为 \(a\)，则 \(a\) 的值是（\quad）

\choices
    {1}
    {2}
    {\( \frac{1}{2} \)}
    {\( \frac{5}{4} \)}
\end{problem}

\begin{problem}
如图，在平面直角坐标系中，\(\Omega\) 是一个与 \(x\) 轴的正半轴，\(y\) 轴的正半轴分别相切于点 \(C\)、\(D\) 的定圆所围成区域 (含边界), \(A\)、\(B\)、\(C\)、\(D\) 是该圆的四等分点. 若点 \(P(x, y)\)、点 \(P'(x', y')\) 满足 \(x \leqslant x'\) 且 \(y \geqslant y'\)，则称 \(P\) 优于 \(P'\). 如果 \(\Omega\) 中的点 \(Q\) 满足: 不存在 \(\Omega\) 中的其它点优于 \(Q\)，那么所有这样的点 \(Q\) 组成的集合是劣弧（\quad）

\bitmapfigure[max width=0.66\linewidth,max totalheight=4.8cm]{img/2008/shanghai_liberal/q15_fig1.png}

\choices
    {弧 \(AB\)}
    {弧 \(BC\)}
    {弧 \(CD\)}
    {弧 \(DA\)}
\end{problem}

\begin{problem}
如图，某住宅小区的平面图呈扇形 \(AOC\). 小区的两个出入口设置在点 \(A\) 及点 \(C\) 处，小区里有两条笔直的小路 \(AD\)、\(DC\)，且拐弯处的转角为 \(120^{\circ}\). 已知某人从 \(C\) 沿 \(CD\) 走到 \(D\) 用了 10 分钟，从 \(D\) 沿 \(DA\) 走到 \(A\) 用了 6 分钟. 若此人步行的速度为每分钟 50 米，求该扇形的半径 \(OA\) 的长. (精确到 1 米)
\end{problem}

\begin{problem}
已知函数 \( f(x) = \sin 2x \)，\( g(x) = \cos \left(2x + \frac{\pi}{6}\right) \)，直线 \( x = t (t \in \R) \) 与函数 \( f(x) \)、\( g(x) \) 的图象分别交于 \( M \)、\( N \) 两点. (1) 当 \( t = \frac{\pi}{4} \) 时，求 \( |MN| \) 的值; (2) 求 \( \left|MN\right| \) 在 \( t \in \left[0, \frac{\pi}{2}\right] \) 时的最大值.
\end{problem}

\begin{problem}
已知函数 \( f(x) = 2^x - \frac{1}{2^{2x}} \). (1) 若 \( f(x) = 2 \)，求 \( x \) 的值； (2) 若 \(2^{x}f(2t) + mf(t) \geqslant 0\) 对于 \(t \in [1,2]\) 恒成立，求实数 \(m\) 的取值范围.
\end{problem}

\begin{problem}
已知双曲线 \(C: \frac{x^2}{3} - y^2 = 1\). (1) 求双曲线 C 的渐近线方程; (2) 已知点 \(M\) 的坐标为 \((0,1)\). 设 \(P\) 是双曲线 \(C\) 上的点，\(Q\) 是点 \(P\) 关于原点的对称点. 记 \(\lambda = \overrightarrow{MP} \cdot \overrightarrow{MQ}\)，求 \(\lambda\) 的取值范围; (3) 已知点 \(D\)、\(E\)、\(M\) 的坐标分别为 \((-2, -1)\)、\((2, -1)\)、\((0, 1)\)，\(P\) 为双曲线 \(C\) 上在第一象限内的点. 记 \(l\) 为经过原点与点 \(P\) 的直线，\(s\) 为 \(\triangle DEM\) 载直线 \(l\) 所得线段的长. 试将 \(s\) 表示为直线 \(l\) 的斜率 \(k\) 的函数. 
\end{problem}

\begin{problem}
已知数列 \(\{a_{n}\}: a_{1} = 1, a_{2} = 2, a_{3} = r, a_{n+3} = a_{n} + 2n\) （\(n\) 是正整数），与数列 \(\{b_{n}\}: b_{1} = 1, b_{2} = 0, b_{3} = -1, b_{4} = 0, b_{n+4} = b_{n}\) （\(n\) 是正整数），记 \(T_{n} = b_{1}a_{1} + b_{2}a_{2} + b_{3}a_{3} + \cdots + b_{n}a_{n}\). (1) 若 \(a_1 + a_2 + a_3 + \cdots + a_{12} = 64\)，求 \(r\) 的值; (2) 求证: 当 n 是正整数时，\( T_{12n} = -4n \) ; (3) 已知 \(r > 0\)，且存在正整数 \(m\)，使得在 \(T_{12m+1}, T_{12m+2}, \cdots, T_{12m+12}\) 中有 4 项为 100，求 \(r\) 的值，并指出哪 4 项为 100.
\end{problem}

